\chapter{The Spectral Theorem and Its Relatives}\label{linear:ch:spectral}

\begin{center}\itshape
``I developed my theory of infinitely many variables from purely mathematical interests, and even called it `spectral analysis' without any presentiment that it would later find application to the actual spectrum of physics.''
\upshape(attributed to David Hilbert, in the 1920s)
\end{center}

A prism separates light into colours. Spectral theory separates an operator into simpler actions. Hilbert borrowed the word \emph{spectrum} while studying integral equations: it names the values $\lambda$ for which $A-\lambda I$ is not invertible. Quantum mechanics later gave the name a physical meaning. Heisenberg, Born, Jordan and Schr\"odinger expressed atomic energy levels through operators; differences between the eigenvalues correspond to spectral lines. Hilbert's reported remark above recalls that unexpected connection.

The finite-dimensional theorem came earlier. Cauchy proved the real symmetric case in 1829: the eigenvalues are real and eigenvectors can be chosen as perpendicular unit vectors. Rayleigh's \emph{Theory of Sound} (1877) approached frequencies by minimising energy ratios. Fischer (1905) and Courant (1920) extended this variational description to every eigenvalue. Weyl used it in 1912 to study the growth of eigenvalues of planar drums and to compare the spectra of $A$, $B$ and $A+B$.

\portrait{cauchy}{Augustin-Louis Cauchy (1789--1857). Illustration from \emph{Das neunzehnte Jahrhundert in Bildnissen}, Berlin, 1901. Public domain, via Wikimedia Commons.}

Normal matrices have orthonormal eigenbases. For Hermitian matrices, their eigenvalues also solve optimisation problems. Rayleigh quotients lead to min-max formulas, interlacing and Weyl's inequalities. Schur and Horn describe the diagonals compatible with a given spectrum. Hoffman--Wielandt bounds the total movement of eigenvalues under perturbation.

\begin{think}
Can a real symmetric matrix have diagonal $(3,0,0)$ and eigenvalues $(2,2,-1)$? Can it have diagonal $(1,1,1)$ and eigenvalues $(3,0,0)$? Both pairs have matching traces. What further obstruction would you look for? Section~\ref{linear:sec:schurhorn} will give a complete test.
\end{think}

Matrices let us examine a spectrum with finitely many coordinates. Chapter~\ref{ch:spectral} will return to infinite dimension and the compact spectral theorem. The finite-dimensional proofs here show which geometric arguments are worth carrying over.

\begin{reminder}[what this chapter assumes]
We use Volume~1, Chapter~7, and Chapter~\ref{ch:hilbert}'s orthonormal bases, orthogonal complements and adjoints. In an orthonormal basis, the adjoint has the conjugate-transpose matrix. Appendices~\ref{linear:ch:eigen} and~\ref{linear:ch:inner} give a separate finite-dimensional account. The one additional fact needed now is Schur's triangularisation; the following argument establishes it before use.
\end{reminder}

A complex operator has an eigenvector because its characteristic polynomial has a root (Volume~1, Theorem~6.4). Choose a unit eigenvector and complete it to an orthonormal basis by Gram--Schmidt. The first column of the matrix has its eigenvalue at the top and zeros below. Apply induction to the lower-right block, using an orthonormal change of basis on those remaining coordinates. This makes the whole matrix upper triangular. Over $\R$, the same induction works whenever the characteristic polynomial splits over $\R$: the lower block's characteristic polynomial is the original one divided by the first eigenvalue factor. This proves the form of Schur's theorem also recorded in Theorem~\ref{linear:thm:schur}.

\section{The spectral theorem}

Unless we specify real scalars, $V$ is a finite-dimensional complex inner product space. The notation $A^*$ always means the adjoint.

\begin{definition}\label{linear:def:normal}
An operator is \emph{normal} when $AA^*=A^*A$. It is \emph{Hermitian}, or \emph{self-adjoint}, when $A^*=A$; over $\R$ this is called symmetric. The conditions $A^*=-A$ and $A^*A=I$ define skew-Hermitian and unitary operators, respectively.
\end{definition}

\begin{theorem}[spectral theorem]\label{linear:thm:spectral}
A finite-dimensional complex operator is normal exactly when it admits an orthonormal eigenbasis. Equivalently, a complex matrix is normal exactly when $A=UDU^*$ for a unitary $U$ and diagonal $D$.
\end{theorem}
\begin{proof}
For $A=UDU^*$, its adjoint is $U\bar DU^*$. Diagonal matrices commute, so $AA^*=UD\bar DU^*=U\bar DDU^*=A^*A$.

Conversely, triangularisation gives a unitary $U$ with $T=U^*AU$ upper triangular. Normality passes to $T$, since $TT^*=U^*AA^*U=U^*A^*AU=T^*T$. The $(1,1)$ entries compare the squared norms of the first row and column. The column contributes only $|t_{11}|^2$, whereas the row contributes $\sum_j|t_{1j}|^2$. Equality forces $t_{1j}=0$ for $j>1$. Thus $T=\operatorname{diag}(t_{11},T')$, where $T'$ remains upper triangular and normal. Induction makes $T$ diagonal.
\end{proof}

\begin{corollary}\label{linear:cor:hermitian}
Hermitian eigenvalues are real, and unitary eigenvalues have modulus $1$. Eigenvectors of a normal operator belonging to different eigenvalues are orthogonal. Every real symmetric matrix has $A=QDQ^{\mathsf T}$ with real orthogonal $Q$ and real diagonal $D$.
\end{corollary}
\begin{proof}
In $A=UDU^*$, the identity $A=A^*$ gives $D=\bar D$, and $A^*A=I$ gives $\bar DD=I$. Distinct eigenspaces are spanned by disjoint parts of the orthonormal eigenbasis and hence are orthogonal. For real symmetric $A$, its eigenvalues are real even when first computed over $\C$, so its characteristic polynomial splits over $\R$. The real triangularisation proved above supplies an orthogonal $Q$ making $Q^{\mathsf T}AQ$ upper triangular. Symmetry makes that triangular matrix diagonal.
\end{proof}

Group eigenvectors by eigenvalue and write $P_\lambda$ for projection onto the corresponding eigenspace. These projections have orthogonal ranges, sum to $I$, and give the \emph{spectral decomposition} $A=\sum_\lambda\lambda P_\lambda$. A function of the operator is then $f(A)=\sum_\lambda f(\lambda)P_\lambda$. This defines $e^A$ and $|A|$; for Hermitian $A$ with nonnegative eigenvalues it also defines $\sqrt A$ (Exercise~\ref{linear:ex:sqrtpsd}). Chapter~\ref{ch:spectral} extends this way of thinking to compact operators.

Call a Hermitian matrix \emph{positive semidefinite}, written $A\succeq0$, when its eigenvalues are nonnegative. Expanding in an eigenbasis shows this is equivalent to $x^*Ax\ge0$ for every $x$. We order all Hermitian spectra decreasingly: $\lambda_1(A)\ge\cdots\ge\lambda_n(A)$.

\section{Rayleigh quotients and the min-max theorem}

For nonzero $x$, define the real-valued Rayleigh quotient $R_A(x)=\langle Ax,x\rangle/\langle x,x\rangle$. In an orthonormal eigenbasis $Au_i=\lambda_iu_i$, write $x=\sum_ic_iu_i$. Then
\begin{equation}\label{linear:eq:rayleigh}
R_A(x)=\frac{\sum_i\lambda_i|c_i|^2}{\sum_i|c_i|^2}.
\end{equation}
It is a weighted average of eigenvalues. Therefore its maximum is $\lambda_1$ and minimum is $\lambda_n$, both attained on eigenvectors:
\[\lambda_1(A)=\max_{x\ne0}R_A(x),\qquad\lambda_n(A)=\min_{x\ne0}R_A(x).\]
On unit vectors the numerator is linear in $A$, so taking its maximum makes $\lambda_1$ convex; in particular $\lambda_1(A+B)\le\lambda_1(A)+\lambda_1(B)$. Also $a_{ii}=R_A(e_i)$ lies between the extreme eigenvalues. This settles the first opening question: an entry $3$ cannot occur on the diagonal when the largest eigenvalue is $2$.

\begin{figure}[ht]
\centering
\fig{linear-ch05-rayleigh}
\caption{The Rayleigh quotient of a $2\times2$ symmetric matrix with eigenvalues $3$ and $1$, on the unit vectors
$u_\theta=(\cos\theta,\sin\theta)$. It is $3\cos^2(\theta-\frac\pi6)+\sin^2(\theta-\frac\pi6)$, maximal on the top eigenvector
and minimal on the other.}
\label{linear:fig:rayleigh}
\end{figure}

To reach an eigenvalue in the middle, first minimise on a subspace, then choose the best subspace.

\begin{theorem}[Courant--Fischer]\label{linear:thm:courantfischer}
For a Hermitian operator on an $n$-dimensional space and $1\le k\le n$, both formulas recover its $k$th largest eigenvalue:
\[\lambda_k(A)=\max_{\dim W=k}\ \min_{x\in W\setminus0}R_A(x)=\min_{\dim W=n-k+1}\ \max_{x\in W\setminus0}R_A(x).\]
\end{theorem}
\begin{proof}
The subspace $W_k=\operatorname{span}(u_1,\dots,u_k)$ has minimum quotient $\lambda_k$. Any other $k$-dimensional $W$ meets $U=\operatorname{span}(u_k,\dots,u_n)$ nontrivially, because their dimensions add to $n+1$. At a nonzero vector of that intersection, formula~\eqref{linear:eq:rayleigh} gives $R_A(x)\le\lambda_k$. Thus $\min_WR_A\le\lambda_k$, proving the max-min identity. Apply that identity to $-A$, whose eigenvalues reverse order and sign, to obtain the min-max identity.
\end{proof}

These formulas compare eigenvalues without matching eigenvectors. Compressing an operator to a smaller subspace gives the first application.

\begin{theorem}[Cauchy interlacing]\label{linear:thm:interlacing}
Let $A$ be Hermitian of size $n$ and $P\in M_{n,m}$ have orthonormal columns. Its compression $B=P^*AP$ satisfies
\[\lambda_k(A)\ge\lambda_k(B)\ge\lambda_{k+n-m}(A)\qquad(1\le k\le m).\]
Principal submatrices are examples. In particular, an $(n-1)$-dimensional principal corner with eigenvalues $\mu_1\ge\cdots\ge\mu_{n-1}$ has
\[\lambda_1\ge\mu_1\ge\lambda_2\ge\mu_2\ge\cdots\ge\mu_{n-1}\ge\lambda_n.\]
\end{theorem}
\begin{proof}
Since $P$ preserves lengths, $R_B(y)=R_A(Py)$. Its images of subspaces are precisely the subspaces of $\operatorname{Im}P$, with dimensions unchanged. Courant--Fischer therefore computes $\lambda_k(B)$ by maximising the minimum over $k$-dimensional subspaces of $\operatorname{Im}P$. Allowing all such subspaces of the ambient space can only increase the answer to $\lambda_k(A)$. For the other bound, minimise the maximum over $(m-k+1)$-dimensional subspaces. Restricting to $\operatorname{Im}P$ can only increase this minimum, whose unrestricted value is $\lambda_{n-(m-k+1)+1}(A)=\lambda_{k+n-m}(A)$.
\end{proof}

\begin{figure}[ht]
\centering
\fig{linear-ch05-interlace}
\caption{The eigenvalues of the top-left $k\times k$ corners of a random symmetric $6\times6$ matrix, for $k=1,\dots,6$.
Each row interlaces the next one: the thin lines show that every eigenvalue of a corner lies between two consecutive
eigenvalues of the next corner. Such triangular arrays are called Gelfand--Tsetlin patterns.}
\label{linear:fig:interlace}
\end{figure}

Appendix~\ref{linear:ch:inner} applies interlacing to successive Jacobi matrices, giving interlacing zeros of orthogonal polynomials (Figure~\ref{linear:fig:legendre}). Chapter~\ref{linear:ch:graphs} applies it to a signed adjacency matrix of the cube and proves the sensitivity conjecture. Conversely, prescribed interlacing sequences can be realised by a matrix and a corner; Exercise~\ref{linear:ex:bordered} constructs one.

\section{Weyl's inequalities}

The spectra of $A$ and $B$ do not determine the spectrum of their sum: the relative positions of eigenvectors matter. They still give useful bounds.

\begin{theorem}[Weyl, 1912]\label{linear:thm:weyl}
If $A$ and $B$ are Hermitian of size $n$ and $i+j-1\le n$, then
\[\lambda_{i+j-1}(A+B)\le\lambda_i(A)+\lambda_j(B).\]
\end{theorem}
\begin{proof}
Let $U$ span the eigenvectors of $A$ from index $i$ onward, and $U'$ those of $B$ from index $j$ onward. With $k=i+j-1$, their intersection has dimension at least $(n-i+1)+(n-j+1)-n=n-k+1$. On $U$, $R_A\le\lambda_i(A)$; on $U'$, $R_B\le\lambda_j(B)$. Choose $W\subseteq U\cap U'$ of dimension $n-k+1$. There $R_{A+B}=R_A+R_B\le\lambda_i(A)+\lambda_j(B)$. The min-max formula gives the result.
\end{proof}

\begin{corollary}[stability of eigenvalues]\label{linear:cor:weylpert}
For Hermitian $A,B$ and every $k$,
\[\lambda_n(B)\le\lambda_k(A+B)-\lambda_k(A)\le\lambda_1(B).\]
Consequently $|\lambda_k(A+B)-\lambda_k(A)|\le\|B\|$, where $\|B\|=\max(|\lambda_1(B)|,|\lambda_n(B)|)$ is the operator norm.
\end{corollary}
\begin{proof}
Set $j=1$ in Weyl's inequality for the upper bound. Apply it again to $(A+B)+(-B)=A$ to obtain the lower bound, since $\lambda_1(-B)=-\lambda_n(B)$. In an eigenbasis of $B$, $\|Bx\|^2=\sum_i\lambda_i(B)^2|c_i|^2\le\max_i|\lambda_i(B)|^2\|x\|^2$, with equality on a corresponding unit eigenvector. This proves the norm formula.
\end{proof}

Hermitian eigenvalues move no farther than the matrix in operator norm. Nonnormal matrices can behave quite differently: Appendix~\ref{linear:ch:eigen} illustrates a perturbation of size $10^{-6}$ moving eigenvalues of a $12\times12$ Jordan block by about $0.3$ (Figure~\ref{linear:fig:perturb}). Exercise~\ref{linear:ex:nonnormal} gives a two-dimensional example with exact formulas.

\begin{remark}
Horn conjectured a complete recursive list of inequalities for the spectra of $A$, $B$ and $A+B$ in 1962. Klyachko and Knutson--Tao proved the conjecture in 1998--1999, using geometric invariant theory and honeycombs; the latter established the needed saturation property. These inequalities also describe tensor-product questions for $GL_n$ and intersections of Schubert varieties. Fulton's survey in the sources explains these connections; none is used in our proofs.
\end{remark}

\section{Diagonals and spectra: Schur and Horn}\label{linear:sec:schurhorn}

For $x\in\R^n$, let $x^\downarrow$ list its coordinates in decreasing order. The relation $x\prec y$, read ``$x$ is majorised by $y$'', means
\[\sum_{i=1}^kx_i^\downarrow\le\sum_{i=1}^ky_i^\downarrow\quad(1\le k\le n),\qquad\sum_{i=1}^nx_i=\sum_{i=1}^ny_i.\]
The partial sums say that $y$ allows larger concentrations of mass. For a probability vector $x$, for example, $(1/3,1/3,1/3)\prec x\prec(1,0,0)$ in dimension three. A \emph{doubly stochastic} matrix has nonnegative entries and every row and column sum equal to $1$. Permutation matrices are examples. A linear expression on such a matrix can always be improved to one on a permutation matrix.

\begin{lemma}\label{linear:lem:birkhoff}
For doubly stochastic $S=(s_{ij})$ and real coefficients $c_{ij}$, some permutation $\sigma$ satisfies $\sum_ic_{i\sigma(i)}\le\sum_{i,j}c_{ij}s_{ij}$.
\end{lemma}
\begin{proof}
Call an entry fractional if it is strictly between $0$ and $1$. With no fractional entries, each row and column has exactly one $1$, so $S$ is a permutation matrix. Otherwise every row or column containing a fractional entry contains at least two, since its sum is the integer $1$. Form a bipartite graph with rows and columns as vertices and fractional entries as edges. Each nonisolated vertex has degree at least two. Following edges without immediately retracing one eventually produces a cycle; bipartiteness makes its length even. Alternately add $\varepsilon$ and subtract $\varepsilon$ around the cycle. Row and column sums do not change. The objective changes by $\varepsilon\delta$ for a fixed real $\delta$. Choose the sign so this is nonpositive, then increase $|\varepsilon|$ until an affected entry first reaches $0$ or $1$. All entries remain in $[0,1]$ and the number of fractional entries decreases. Finitely many repetitions produce the desired permutation.
\end{proof}

This cycle argument underlies Birkhoff's theorem that doubly stochastic matrices are convex combinations of permutation matrices. The lemma alone is enough for our eigenvalue estimates.

\begin{theorem}[Schur, 1923]\label{linear:thm:schurmaj}
For a Hermitian matrix, its diagonal is majorised by its eigenvalue vector: $(a_{11},\dots,a_{nn})\prec(\lambda_1,\dots,\lambda_n)$.
\end{theorem}
\begin{proof}
From $A=UDU^*$, each $a_{ii}=\sum_j|u_{ij}|^2\lambda_j$. The matrix $(|u_{ij}|^2)$ is doubly stochastic, since rows and columns of $U$ are unit vectors. The total sums agree by the trace identity. For the $k$ largest diagonal entries, let $I$ be their index set and set $t_j=\sum_{i\in I}|u_{ij}|^2$. Then $0\le t_j\le1$ and $\sum_jt_j=k$. Since the eigenvalues are decreasing,
\[\begin{aligned}
\sum_{j\le k}\lambda_j-\sum_{i\in I}a_{ii}
&=\sum_{j\le k}(1-t_j)\lambda_j-\sum_{j>k}t_j\lambda_j\\
&\ge\lambda_k\left(\sum_{j\le k}(1-t_j)-\sum_{j>k}t_j\right)=0.
\end{aligned}\]
These are exactly the majorisation inequalities.
\end{proof}

Changing the orthonormal basis shows that $\sum_{i\le k}\langle Ax_i,x_i\rangle\le\lambda_1+\cdots+\lambda_k$ for every orthonormal family of size $k$. The top eigenvectors attain equality. This is Ky Fan's maximum principle, developed in Exercise~\ref{linear:ex:kyfan}. Majorisation also gives a sufficient condition for a diagonal, not just a necessary one.

\begin{theorem}[Horn, 1954]\label{linear:thm:horn}
Whenever $d\prec\lambda$ in $\R^n$, a real symmetric matrix exists with diagonal $d$ and eigenvalues $\lambda$.
\end{theorem}
\begin{proof}
We construct an orthogonal conjugate of $\Lambda=\operatorname{diag}(\lambda)$ with the prescribed diagonal. Permutations let us assume $d$ and $\lambda$ are decreasing. For $n=1$, equality of totals gives the result. Suppose $n\ge2$ and use induction.

Majorisation gives $\lambda_1\ge d_1$, while $d_1\ge n^{-1}\sum_i d_i=n^{-1}\sum_i\lambda_i\ge\lambda_n$. Choose $k\le n-1$ with $\lambda_k\ge d_1\ge\lambda_{k+1}$. Some $0\le c\le1$ satisfies $c^2\lambda_k+(1-c^2)\lambda_{k+1}=d_1$. Put $s=\sqrt{1-c^2}$ and rotate these two coordinates:
\[\begin{pmatrix}c&-s\\s&c\end{pmatrix}
\begin{pmatrix}\lambda_k&0\\0&\lambda_{k+1}\end{pmatrix}
\begin{pmatrix}c&s\\-s&c\end{pmatrix}
=\begin{pmatrix}d_1&cs(\lambda_k-\lambda_{k+1})\\cs(\lambda_k-\lambda_{k+1})&\lambda_k+\lambda_{k+1}-d_1\end{pmatrix}.\]
Move the first of these coordinates to the front. The resulting matrix has form $M=\bigl(\begin{smallmatrix}d_1&r^{\mathsf T}\\r&\Lambda'\end{smallmatrix}\bigr)$, with $\Lambda'$ diagonal and entries
\[\lambda'=(\lambda_1,\dots,\lambda_{k-1},\lambda_k+\lambda_{k+1}-d_1,\lambda_{k+2},\dots,\lambda_n).\]
The inserted entry lies between $\lambda_{k+1}$ and $\lambda_k$, so this list is decreasing. We need $(d_2,\dots,d_n)\prec\lambda'$. The totals agree. For $m<k$, the first $m$ entries of $\lambda'$ sum to $\sum_{i\le m}\lambda_i\ge m\lambda_k\ge md_1\ge\sum_{i=2}^{m+1}d_i$. For $m\ge k$, they sum to $\sum_{i\le m+1}\lambda_i-d_1\ge\sum_{i=2}^{m+1}d_i$. Induction supplies $Q'$ with diagonal $(d_2,\dots,d_n)$ in $Q'\Lambda'Q'^{\mathsf T}$. Conjugating $M$ by $\operatorname{diag}(1,Q')$ leaves its first diagonal entry fixed and installs the remaining ones. It stays orthogonally similar to $\Lambda$, so its spectrum is unchanged.
\end{proof}

The all-ones matrix answers the second opening question: its diagonal is $(1,1,1)$ and its spectrum is $(3,0,0)$. Horn's proof constructs all admissible diagonals with successive plane rotations, fixing one entry at a time. Exercise~\ref{linear:ex:hornsmall} gives a small calculation.

\section{How far eigenvalues move}

The operator norm bounds each Hermitian eigenvalue separately. To bound their combined movement, use the Frobenius norm $\|M\|_F=(\sum_{i,j}|m_{ij}|^2)^{1/2}=\sqrt{\operatorname{tr}M^*M}$. Unitary changes of basis preserve it. The resulting theorem applies even to normal matrices with complex spectra.

\begin{theorem}[Hoffman--Wielandt, 1953]\label{linear:thm:hoffmanwielandt}
For normal matrices $A,B$ with eigenvalues $\alpha_1,\dots,\alpha_n$ and $\beta_1,\dots,\beta_n$, some permutation $\sigma$ satisfies
\[\sum_{i=1}^n|\alpha_i-\beta_{\sigma(i)}|^2\le\|A-B\|_F^2.\]
For Hermitian $A,B$, decreasing order is an admissible pairing: $\sum_i(\lambda_i(A)-\lambda_i(B))^2\le\|A-B\|_F^2$.
\end{theorem}
\begin{proof}
Write $A=UD_\alpha U^*$, $B=VD_\beta V^*$ and $W=U^*V$. Unitary invariance gives
\[\|A-B\|_F^2=\|D_\alpha W-WD_\beta\|_F^2=\sum_{i,j}|\alpha_i-\beta_j|^2|w_{ij}|^2.\]
The last identity follows entry by entry. Since $(|w_{ij}|^2)$ is doubly stochastic, Lemma~\ref{linear:lem:birkhoff} with $c_{ij}=|\alpha_i-\beta_j|^2$ gives the required permutation. For real spectra, an inverted pairing with $\alpha_i>\alpha_l$ and $\beta_{\sigma(i)}<\beta_{\sigma(l)}$ can be improved by swapping the two partners. The change is
\[-2(\alpha_i-\alpha_l)(\beta_{\sigma(l)}-\beta_{\sigma(i)})<0.\]
Thus a minimum can be chosen to match both decreasing lists.
\end{proof}

\begin{beauty}
An eigenvalue starts as a root of a polynomial, yet its variational description turns comparison into geometry. Min-max gives interlacing and stability. Majorisation describes exactly which diagonals can accompany a fixed Hermitian spectrum. These diagonals form a convex polytope, the permutohedron; Rado's theorem identifies it with the convex hull of the permuted spectrum. The two pictures agree: rotating an eigenbasis redistributes the diagonal entries without changing their total.
\end{beauty}

\section{From eigenvalues to singular values}

Rayleigh quotients rule out diagonal $(3,0,0)$ for spectrum $(2,2,-1)$. Schur--Horn accepts diagonal $(1,1,1)$ for spectrum $(3,0,0)$. Both tests rely on the orthonormal eigenbasis supplied by the spectral theorem. Min-max and doubly stochastic matrices then quantify comparisons: corners interlace, Hermitian perturbations have controlled eigenvalue shifts, and Hoffman--Wielandt pairs entire normal spectra.

A rectangular matrix has no eigenvalue equation on a single underlying space. Chapter~\ref{linear:ch:svd} replaces eigenvalues by singular values and uses them to find the best approximation by a matrix of smaller rank.

\begin{figure}[ht]
\centering\fig{ch03-rw-diagonals}
\caption{For spectrum $(3,2,1)$, Schur--Horn allows exactly the triples with sum $6$ and $1\le d_i\le3$. In that plane these inequalities cut out a hexagon. Its centre is $(2,2,2)$; the six corners are permutations of the spectrum.}\label{fig:v8c03diagonals}
\end{figure}
\begin{figure}[ht]
\centering\fig{ch03-rw-shifts}
\caption{Adding $\varepsilon I$ to $\operatorname{diag}(3,1)$ shifts each eigenvalue by $\varepsilon$. Both Weyl's individual bound and Hoffman--Wielandt's sum of squared shifts are equalities.}\label{fig:v8c03shifts}
\end{figure}

\section{Exercises}

\begin{exercise}[$\star$]\label{linear:ex:normaliff}
Prove that $A$ is normal exactly when $\|Ax\|=\|A^*x\|$ for every $x$. Apply this criterion to show that $Ax=\lambda x$ for normal $A$ implies $A^*x=\bar\lambda x$.
\end{exercise}
\begin{exercise}[$\star$]\label{linear:ex:hornsmall}
Use the rotations in Horn's proof to construct a real symmetric matrix with diagonal $(2,1,0)$ and spectrum $(3,0,0)$. Decide whether diagonal $(2,2,-1)$ is possible for that same spectrum.
\end{exercise}
\begin{exercise}[$\star$]\label{linear:ex:pathgraph}
The path adjacency matrix $A_n$ has ones immediately above and below the diagonal and zeros elsewhere. Prove that $v_k=(\sin(jk\pi/(n+1)))_{j=1}^n$ is an eigenvector with eigenvalue $2\cos(k\pi/(n+1))$. Verify interlacing between $A_n$ and its corner $A_{n-1}$.
\end{exercise}
\begin{exercise}[$\star\star$]\label{linear:ex:degree}
For a finite graph's adjacency matrix $A$, with average degree $\bar d$ and maximum degree $\Delta$, prove $\bar d\le\lambda_1(A)\le\Delta$. Show that the first inequality is an equality exactly for regular graphs.
\end{exercise}
\begin{exercise}[$\star$]\label{linear:ex:nonnormal}
Take $A=\bigl(\begin{smallmatrix}0&1\\0&0\end{smallmatrix}\bigr)$ and $E=\bigl(\begin{smallmatrix}0&0\\\varepsilon&0\end{smallmatrix}\bigr)$ with $\varepsilon>0$. Compute the spectrum of $A+E$ and show that the bounds of Corollary~\ref{linear:cor:weylpert} and Theorem~\ref{linear:thm:hoffmanwielandt} fail for these matrices when $\varepsilon<1$.
\end{exercise}
\begin{exercise}[$\star\star$]\label{linear:ex:rankone}
If $A$ is Hermitian and $v$ is a vector, prove the rank-one interlacing inequalities
\[\lambda_1(A+vv^*)\ge\lambda_1(A)\ge\lambda_2(A+vv^*)\ge\cdots\ge\lambda_n(A+vv^*)\ge\lambda_n(A).\]
\end{exercise}
\begin{exercise}[$\star\star$]\label{linear:ex:kyfan}
\emph{Ky Fan's principle (1949).} Prove $\lambda_1(A)+\cdots+\lambda_k(A)=\max\sum_{i=1}^k\langle Ax_i,x_i\rangle$, where the maximum ranges over orthonormal families of size $k$. Deduce convexity of the sum of the top $k$ eigenvalues and the inequality $\sum_{i\le k}\lambda_i(A+B)\le\sum_{i\le k}\lambda_i(A)+\sum_{i\le k}\lambda_i(B)$.
\end{exercise}
\begin{exercise}[$\star\star$]\label{linear:ex:sqrtpsd}
For positive semidefinite $A$, prove existence and uniqueness of a positive semidefinite square root. Show that this square root equals a polynomial in $A$.
\end{exercise}
\begin{exercise}[$\star\star$]\label{linear:ex:simultaneous}
Prove that two commuting Hermitian matrices share an orthonormal eigenbasis. Determine whether commuting normal matrices have the same property.
\end{exercise}
\begin{exercise}[$\star\star$]\label{linear:ex:realnormal}
Prove that every real normal $A$ has an orthogonal change of basis making it block diagonal, with blocks $(a)$ or $\bigl(\begin{smallmatrix}a&-b\\b&a\end{smallmatrix}\bigr)$. Interpret this form for a rotation of $\R^3$.
\end{exercise}
\begin{exercise}[$\star\star\star$]\label{linear:ex:bordered}
Given strict interlacing $\lambda_1>\mu_1>\lambda_2>\cdots>\mu_{n-1}>\lambda_n$, construct real $c$ and $b_i$ so that $\bigl(\begin{smallmatrix}D&b\\b^{\mathsf T}&c\end{smallmatrix}\bigr)$, with $D=\operatorname{diag}(\mu_1,\dots,\mu_{n-1})$, has spectrum $(\lambda_1,\dots,\lambda_n)$. This proves a converse to Cauchy interlacing.
\end{exercise}
\begin{exercise}[$\star\star\star$]\label{linear:ex:vonneumann}
\emph{Von Neumann's trace inequality for Hermitian matrices.} Prove $\operatorname{tr}(AB)\le\sum_i\lambda_i(A)\lambda_i(B)$ for Hermitian $A,B$, and characterise equality.
\end{exercise}

\begin{exercise}[$\star$]\label{ex:v8c03a}
For $A=\bigl(\begin{smallmatrix}2&1\\1&2\end{smallmatrix}\bigr)$, find an orthonormal eigenbasis and its two spectral projections. Use them to compute $A^m$ for every integer $m\ge0$, then give $A^3$ explicitly.
\end{exercise}
\begin{exercise}[$\star\star$]\label{ex:v8c03b}
Compress $D=\operatorname{diag}(4,2,-1)$ to the subspace with orthonormal basis $(e_1+e_3)/\sqrt2,e_2$. Compute the compression matrix and its ordered eigenvalues. Check each interlacing bound against the spectrum of $D$.
\end{exercise}
\begin{exercise}[$\star$]\label{ex:v8c03c}
For nonnegative $a,b,c$ with $a+b+c=1$, construct a positive semidefinite matrix with diagonal $(a,b,c)$ and spectrum $(1,0,0)$. Give its entries for $(a,b,c)=(1/2,1/3,1/6)$. Why could no negative diagonal entry occur for this spectrum?
\end{exercise}
\begin{exercise}[$\star\star$]\label{ex:v8c03d}
Let $A=\operatorname{diag}(3,1)$ and $\varepsilon>0$. Show that $A+\varepsilon I$ attains equality in both Weyl's perturbation bound and Hoffman--Wielandt. Next replace the perturbation by $\bigl(\begin{smallmatrix}0&\varepsilon\\\varepsilon&0\end{smallmatrix}\bigr)$. Compute the two eigenvalue shifts and prove that their common absolute value is at most $\varepsilon^2/2$.
\end{exercise}

\section{Hints}

\begin{hint}{linear:ex:normaliff}
Apply the quadratic form of $A^*A-AA^*$ to every vector, then to an eigenbasis of that Hermitian matrix. Use the result for $A-\lambda I$.
\end{hint}
\begin{hint}{linear:ex:hornsmall}
The first rotation needs $3c^2=2$. For the proposed second diagonal, compare the sums of its largest two entries.
\end{hint}
\begin{hint}{linear:ex:pathgraph}
Use $\sin((j-1)\theta)+\sin((j+1)\theta)=2\cos\theta\sin(j\theta)$, including the zero boundary values.
\end{hint}
\begin{hint}{linear:ex:degree}
Test the Rayleigh quotient on the all-ones vector. For the upper bound, inspect an eigenvector at a coordinate of maximum absolute value.
\end{hint}
\begin{hint}{linear:ex:nonnormal}
Expand the characteristic polynomial to get $X^2-\varepsilon$.
\end{hint}
\begin{hint}{linear:ex:rankone}
The perturbation is positive semidefinite, with at most one nonzero eigenvalue. Combine the stability corollary with Weyl's inequality at $j=2$.
\end{hint}
\begin{hint}{linear:ex:kyfan}
Start an orthonormal basis with the given family and apply Schur's majorisation theorem to the resulting diagonal.
\end{hint}
\begin{hint}{linear:ex:sqrtpsd}
Interpolate the square-root function on the finite spectrum. A second square root must commute with $A$ and preserve its eigenspaces.
\end{hint}
\begin{hint}{linear:ex:simultaneous}
Use the eigenspaces of $A$ as invariant subspaces for $B$. In the normal case, interpolation makes $B^*$ a polynomial in $B$.
\end{hint}
\begin{hint}{linear:ex:realnormal}
Pair the conjugate eigenvectors $x+iy$ and $x-iy$. Their orthogonality gives $x\perp y$ and $\|x\|=\|y\|$.
\end{hint}
\begin{hint}{linear:ex:bordered}
Compare $\det(XI-M)/\prod_i(X-\mu_i)$ with the partial fractions of $\prod_j(X-\lambda_j)/\prod_i(X-\mu_i)$. Interlacing fixes the signs of the residues.
\end{hint}
\begin{hint}{linear:ex:vonneumann}
Expand the trace in two eigenbases and use the doubly stochastic lemma and rearrangement. For equality, vary $B$ by plane rotations in an eigenbasis of $A$; equality must remain a maximum.
\end{hint}
\begin{hint}{ex:v8c03a}
Use the unit vectors $(1,1)/\sqrt2$ and $(1,-1)/\sqrt2$. Their projection matrices add to $I$.
\end{hint}
\begin{hint}{ex:v8c03b}
The quotient on the mixed first vector is the average of $4$ and $-1$.
\end{hint}
\begin{hint}{ex:v8c03c}
Take an outer product of a unit vector whose squared coordinates are $a,b,c$.
\end{hint}
\begin{hint}{ex:v8c03d}
For the second perturbation, the characteristic equation is $(X-2)^2=1+\varepsilon^2$. Rationalise $\sqrt{1+\varepsilon^2}-1$.
\end{hint}

\section{Solutions}

\begin{solution}{linear:ex:normaliff}
\idea The difference of the two squared lengths is the quadratic form of a Hermitian matrix. Such a form determines the matrix.
\writeup Put $H=A^*A-AA^*$. Then $\|Ax\|^2-\|A^*x\|^2=\langle Hx,x\rangle$. Normality makes this zero. Conversely, if it vanishes for all $x$, test on an eigenbasis of $H$ to see that every eigenvalue is zero, hence $H=0$. A normal $A$ also makes $A-\lambda I$ normal. Its norm identity applied to $Ax=\lambda x$ gives $\|(A^*-\bar\lambda I)x\|=0$, proving the second assertion.
\end{solution}
\begin{solution}{linear:ex:hornsmall}
\idea One plane rotation already installs the two nonzero diagonal entries in this rank-one case.
\writeup Take $c=\sqrt{2/3}$ and $s=\sqrt{1/3}$. Rotating $\operatorname{diag}(3,0)$ gives $\bigl(\begin{smallmatrix}2&\sqrt2\\\sqrt2&1\end{smallmatrix}\bigr)$. Appending a zero row and column gives
\[M=\begin{pmatrix}2&\sqrt2&0\\\sqrt2&1&0\\0&0&0\end{pmatrix},\]
which has diagonal $(2,1,0)$ and spectrum $(3,0,0)$. For diagonal $(2,2,-1)$, the top two entries sum to $4$, exceeding the top-two eigenvalue sum $3$. Schur's necessary condition rules it out.
\end{solution}
\begin{solution}{linear:ex:pathgraph}
\idea A sampled sine solves the second-order recurrence for an eigenvector of a path.
\writeup With $\theta=k\pi/(n+1)$, each coordinate satisfies $(A_nv_k)_j=\sin((j-1)\theta)+\sin((j+1)\theta)=2\cos\theta\sin(j\theta)$. The boundary terms are zero because $\sin0=\sin((n+1)\theta)=0$. The $n$ eigenvalues $2\cos(k\pi/(n+1))$ are distinct and therefore exhaust the spectrum. For $1\le k\le n-1$, $k/(n+1)\le k/n\le(k+1)/(n+1)$. Cosine decreases on $[0,\pi]$, giving the two required interlacing inequalities.
\end{solution}
\begin{solution}{linear:ex:degree}
\idea One test vector gives the lower bound; one coordinate of a genuine eigenvector gives the upper bound.
\writeup The all-ones vector has quotient $R_A(\mathbf1)=2|E|/n=\bar d$. Thus $\lambda_1\ge\bar d$. Equality in the weighted-average formula~\eqref{linear:eq:rayleigh} forces $\mathbf1$ to lie in the top eigenspace, meaning all row sums, hence all degrees, are equal. A regular graph attains that equality. For the upper bound choose $Ax=\lambda_1x$ and a nonzero coordinate with maximum $|x_i|$. Since $\lambda_1\ge\bar d\ge0$,
\[\lambda_1|x_i|=\left|\sum_{j\sim i}x_j\right|\le\deg(i)|x_i|\le\Delta|x_i|.\]
Divide by $|x_i|$.
\end{solution}
\begin{solution}{linear:ex:nonnormal}
\idea An exact two-dimensional calculation shows why a normality hypothesis matters.
\writeup The characteristic polynomial is $X^2-\varepsilon$, so $A+E$ has eigenvalues $\pm\sqrt\varepsilon$, while $A$ has two zero eigenvalues. Both norms of $E$ are $\varepsilon$. For $0<\varepsilon<1$, each shift has size $\sqrt\varepsilon>\varepsilon$, and the sum of squared shifts is $2\varepsilon>\varepsilon^2=\|E\|_F^2$. These violate the proposed bounds when their matrix hypotheses are omitted.
\end{solution}
\begin{solution}{linear:ex:rankone}
\idea Positivity pushes all eigenvalues upward. Only one positive direction is available, limiting how far the lower eigenvalues can rise.
\writeup Since $B=vv^*\succeq0$, the stability corollary gives $\lambda_k(A+B)\ge\lambda_k(A)$. For $n\ge2$, $B$ has eigenvalues $\|v\|^2,0,\dots,0$. Weyl with $i=k-1$ and $j=2$ yields $\lambda_k(A+B)\le\lambda_{k-1}(A)$ for $k\ge2$. Together these are the asserted interlacing chain. In dimension one, only the positivity inequality remains.
\end{solution}
\begin{solution}{linear:ex:kyfan}
\idea Schur bounds the sum of the first $k$ diagonal entries in any orthonormal basis. An eigenbasis attains the bound.
\writeup Extend the chosen family to an orthonormal basis. The matrix of $A$ in that basis has diagonal entries $\langle Ax_i,x_i\rangle$ and unchanged eigenvalues. Any $k$ diagonal entries sum to at most the top $k$ diagonal entries, whose sum is at most $\lambda_1+\cdots+\lambda_k$. Equality is attained by the top eigenvectors. A maximum of linear functions of $A$ is convex. Finally, the maximum of a sum is at most the sum of the separate maxima, proving the bound for $A+B$.
\end{solution}
\begin{solution}{linear:ex:sqrtpsd}
\idea The spectral decomposition defines the candidate. Any other positive square root must have the same value on each eigenspace.
\writeup Define $S=\sum_\lambda\sqrt\lambda P_\lambda$. It is positive semidefinite and $S^2=A$. Interpolating $p(\lambda)=\sqrt\lambda$ at the finitely many distinct eigenvalues gives $p(A)=S$. If $R\succeq0$ and $R^2=A$, it commutes with $A$ and preserves each eigenspace $E_\lambda$. The restriction of $R$ there is Hermitian, with nonnegative eigenvalues $\mu$ satisfying $\mu^2=\lambda$. Every such eigenvalue is $\sqrt\lambda$, so the restriction is $\sqrt\lambda I$. Thus $R=S$.
\end{solution}
\begin{solution}{linear:ex:simultaneous}
\idea Commutation prevents $B$ from mixing different eigenspaces of $A$. Diagonalise within each of those spaces.
\writeup If $Ax=\lambda x$, then $A(Bx)=B(Ax)=\lambda Bx$, so $B$ preserves $E_\lambda(A)$. For Hermitian $B$, each restriction is Hermitian. Choose its orthonormal eigenbasis on each $E_\lambda$ and unite them; the eigenspaces of $A$ are mutually orthogonal. For normal $B$, use interpolation on its spectrum to write $B^*=p(B)$. It too commutes with $A$, so both $B$ and $B^*$ preserve each $E_\lambda$. The restrictions are therefore normal with adjoint $B^*|_{E_\lambda}$, and the same construction works.
\end{solution}
\begin{solution}{linear:ex:realnormal}
\idea A conjugate pair of complex eigendirections becomes a real invariant plane.
\writeup For a nonreal eigenvalue $a+ib$, choose $x+iy$ with real $x,y$ and $b\ne0$. The eigenvector equation gives $Ax=ax-by$ and $Ay=bx+ay$. Since $A$ is real, $x-iy$ is an eigenvector for $a-ib$ and is orthogonal to $x+iy$. Expanding gives $\|x\|^2-\|y\|^2+2i\langle x,y\rangle=0$, so $x\perp y$ and their norms are equal and nonzero. In the normalised basis $(y,x)$, the invariant plane has matrix $\bigl(\begin{smallmatrix}a&-b\\b&a\end{smallmatrix}\bigr)$. Exercise~\ref{linear:ex:normaliff} shows that $A^*$ also preserves this plane. Its orthogonal complement is therefore invariant for both $A$ and $A^*$, with normal restriction. Real eigenvalues similarly yield real invariant lines. Induction gives the block form. For a rotation in $\R^3$, orthogonality makes each real block $\pm1$ and each plane block a rotation. Their determinant product is $1$, so at least one real block is $1$. Its eigenvector is the rotation axis.
\end{solution}
\begin{solution}{linear:ex:bordered}
\idea Interlacing makes the residues have exactly the signs needed to be negatives of squares.
\writeup Put $f(X)=\prod_j(X-\lambda_j)$ and $g(X)=\prod_i(X-\mu_i)$. Expanding the bordered determinant gives
\[\det(XI-M)=g(X)\left(X-c-\sum_i\frac{b_i^2}{X-\mu_i}\right).\]
Partial fractions give $f/g=X-c_0-\sum_i r_i/(X-\mu_i)$, where $c_0=\sum_j\lambda_j-\sum_i\mu_i$ and $r_i=-f(\mu_i)/g'(\mu_i)$. Strict interlacing makes $f(\mu_i)$ have $i$ negative factors and $g'(\mu_i)$ have $i-1$. Hence $r_i>0$. Choosing $c=c_0$ and $b_i=\sqrt{r_i}$ makes the two polynomials equal, which gives the prescribed spectrum.
\end{solution}
\begin{solution}{linear:ex:vonneumann}
\idea Two eigenbases express the trace as a doubly stochastic average. Rearrangement maximises the pairing; equality forces the two matrices to share compatible eigendirections.
\writeup With $A=UD_\alpha U^*$, $B=VD_\beta V^*$ and $W=U^*V$, we have $\operatorname{tr}(AB)=\sum_{i,j}\alpha_i\beta_j|w_{ij}|^2$. Apply Lemma~\ref{linear:lem:birkhoff} to $-\alpha_i\beta_j$ to find a permutation with $\operatorname{tr}(AB)\le\sum_i\alpha_i\beta_{\sigma(i)}$. Removing an inverted pair increases this sum by $(\alpha_i-\alpha_l)(\beta_{\sigma(l)}-\beta_{\sigma(i)})\ge0$. Thus decreasing lists give its maximum.

A common orthonormal basis ordering both spectra decreasingly gives equality. Conversely, assume equality and work in a basis where $A=\operatorname{diag}(\alpha)$ is decreasing. Among all unitary conjugates of $B$, the present trace is maximal. Conjugate $B$ by a real rotation of coordinates $i,j$ and differentiate the trace at angle zero. Maximality gives $(\alpha_i-\alpha_j)\operatorname{Re}b_{ij}=0$. Multiplying coordinate $i$ by $i$ first preserves the trace and replaces this real part by $-\operatorname{Im}b_{ij}$, so $(\alpha_i-\alpha_j)\operatorname{Im}b_{ij}=0$ too. Hence $B$ preserves every eigenspace of $A$. Diagonalise $B$ within those spaces. If two resulting pairs satisfy $\alpha_i>\alpha_j$ but $\beta_i<\beta_j$, swapping them strictly increases the trace, contradicting maximality. Ordering within tied blocks therefore gives a common eigenbasis ordering both spectra decreasingly. This proves the full equality characterisation.
\end{solution}

\begin{solution}{ex:v8c03a}
\idea Once the eigenspaces are known, powers act separately on the two perpendicular lines.
\writeup The eigenvectors $u_+=(1,1)/\sqrt2$ and $u_-=(1,-1)/\sqrt2$ have eigenvalues $3$ and $1$. Their projections are $P_+=\frac12\bigl(\begin{smallmatrix}1&1\\1&1\end{smallmatrix}\bigr)$ and $P_-=\frac12\bigl(\begin{smallmatrix}1&-1\\-1&1\end{smallmatrix}\bigr)$. Since $P_+P_-=0$ and $P_++P_-=I$, $A^m=3^mP_++P_-$ for every $m\ge0$. Thus $A^3=\bigl(\begin{smallmatrix}14&13\\13&14\end{smallmatrix}\bigr)$.
\end{solution}
\begin{solution}{ex:v8c03b}
\idea Compression keeps the quadratic form on the chosen plane. The mixed vector averages two diagonal entries.
\writeup For $u=(e_1+e_3)/\sqrt2$ and $v=e_2$, $\langle Du,u\rangle=(4-1)/2=3/2$, $\langle Dv,v\rangle=2$, and the cross terms vanish. The compression is $\operatorname{diag}(3/2,2)$, with ordered spectrum $(2,3/2)$. Interlacing reads $4\ge2\ge2$ and $2\ge3/2\ge-1$, as required.
\end{solution}
\begin{solution}{ex:v8c03c}
\idea A unit vector's outer product is projection onto its span, with prescribed squared coordinates on the diagonal.
\writeup Set $v=(\sqrt a,\sqrt b,\sqrt c)$ and $A=vv^{\mathsf T}$. Then $\|v\|=1$, $Av=v$, and $A$ vanishes on $v^\perp$, proving the spectrum $(1,0,0)$ and positivity. Its diagonal is $(a,b,c)$. For the specified triple,
\[A=\begin{pmatrix}1/2&1/\sqrt6&1/\sqrt{12}\\1/\sqrt6&1/3&1/\sqrt{18}\\1/\sqrt{12}&1/\sqrt{18}&1/6\end{pmatrix}.\]
Any matrix with this Hermitian spectrum is positive semidefinite, so $a_{ii}=e_i^*Ae_i\ge0$. Negative diagonal entries are impossible.
\end{solution}
\begin{solution}{ex:v8c03d}
\idea A scalar perturbation shifts all eigendirections equally. The off-diagonal perturbation has an exact quadratic formula.
\writeup Adding $\varepsilon I$ gives eigenvalues $3+\varepsilon,1+\varepsilon$. Each shift equals the perturbation norm $\varepsilon$, and their squared sum is $2\varepsilon^2=\|\varepsilon I\|_F^2$. For the second perturbation, the eigenvalues solve $(X-2)^2=1+\varepsilon^2$, so they are $2\pm\sqrt{1+\varepsilon^2}$. Their shifts from $3,1$ are $\delta,-\delta$, where
\[\delta=\sqrt{1+\varepsilon^2}-1=\frac{\varepsilon^2}{\sqrt{1+\varepsilon^2}+1}\le\frac{\varepsilon^2}{2}.\]
This proves the claimed smaller movement without a numerical approximation.
\end{solution}

\section*{Sources and further reading}
\begin{itemize}
\item A.-L.~Cauchy, \emph{Sur l'\'equation \`a l'aide de laquelle on d\'etermine les in\'egalit\'es s\'eculaires des mouvements des plan\`etes}, Exercices de math\'ematiques 4 (1829), treats real eigenvalues and interlacing.
\item J.~W.~Strutt, Lord Rayleigh, \emph{The Theory of Sound}, Macmillan, 1877.
\item D.~Hilbert, \emph{Grundz\"uge einer allgemeinen Theorie der linearen Integralgleichungen} (Erste Mitteilung), G\"ottinger Nachrichten (1904), 49--91; L.~A.~Steen, \emph{Highlights in the history of spectral theory}, American Mathematical Monthly 80 (1973), 359--381, gives the historical account and attributed quotation.
\item E.~Fischer, \emph{\"Uber quadratische Formen mit reellen Koeffizienten}, Monatshefte f\"ur Mathematik und Physik 16 (1905), 234--249; R.~Courant, \emph{\"Uber die Eigenwerte bei den Differentialgleichungen der mathematischen Physik}, Mathematische Zeitschrift 7 (1920), 1--57.
\item H.~Weyl, \emph{Das asymptotische Verteilungsgesetz der Eigenwerte linearer partieller Differentialgleichungen}, Mathematische Annalen 71 (1912), 441--479.
\item I.~Schur, \emph{\"Uber eine Klasse von Mittelbildungen mit Anwendungen auf die Determinantentheorie}, Sitzungsberichte der Berliner Mathematischen Gesellschaft 22 (1923), 9--20; A.~Horn, \emph{Doubly stochastic matrices and the diagonal of a rotation matrix}, American Journal of Mathematics 76 (1954), 620--630; N.~N.~Chan and K.-H.~Li, \emph{Diagonal elements and eigenvalues of a real symmetric matrix}, Journal of Mathematical Analysis and Applications 91 (1983), 562--566, presents the rotation proof used here.
\item A.~J.~Hoffman and H.~W.~Wielandt, \emph{The variation of the spectrum of a normal matrix}, Duke Mathematical Journal 20 (1953), 37--39; G.~Birkhoff, \emph{Tres observaciones sobre el \'algebra lineal}, Revista Universidad Nacional de Tucum\'an, Serie A, 5 (1946), 147--151.
\item K.~Fan, \emph{On a theorem of Weyl concerning eigenvalues of linear transformations I}, Proceedings of the National Academy of Sciences 35 (1949), 652--655.
\item W.~Fulton, \emph{Eigenvalues, invariant factors, highest weights, and Schubert calculus}, Bulletin of the AMS 37 (2000), 209--249; A.~Knutson and T.~Tao, \emph{Honeycombs and sums of Hermitian matrices}, Notices of the AMS 48 (2001), 175--186.
\item R.~Bhatia, \emph{Matrix Analysis}, Springer, 1997, Chapters II--III and VI; R.~A.~Horn and C.~R.~Johnson, \emph{Matrix Analysis}, 2nd ed., Cambridge University Press, 2013, Chapter~4.
\end{itemize}
